\section{Discussion and limitations}
\label{sec:discussion}

The combined evidence supports a deployment workflow within this benchmark:
retain an explicit belief over persistent physical state, select the
conservatism margin on target-task development lifetimes using a buffered
tolerance, freeze the choice, and evaluate it once on fresh lifetimes. A margin
can improve average utility on a new task yet still miss a prespecified
trip-rate tolerance, so positive transfer of reward is not sufficient evidence
for transferring a numerical operating threshold.

The utility improvement is not free. In Pusher, avoided trip cost outweighs a
decrease in immediate base-task return, and fixed-trajectory reweighting places
the mean break-even trip penalty near 40 under the original throughput bonus.
This accounting sensitivity is not policy retraining and does not establish a
benefit when trips carry little cost.

Both tasks use MuJoCo and share one phenomenological thermal law. The study has
no physical-unit calibration, hardware telemetry, formal safety proof, or
probability-calibrated OOD guarantee. It tests one selected recurrent baseline
rather than exhausting recurrent, meta-RL, system-identification, or online
adaptation methods. These boundaries prevent the result from supporting a
general thermal-management or universal algorithm-ranking claim.

\section{Broader impact}

Persistent-health benchmarks may help reveal controllers that trade short-term
performance for lower accumulated risk. The same abstractions could also be
misused to overstate safety when simulated health signals are poorly calibrated
or when an aggregate trip tolerance hides rare severe outcomes. We mitigate
this risk by preserving failed gates, reporting distributional cautions,
avoiding physical-validity claims, and releasing the exact simulated protocol
needed to inspect the result. The computational cost and environmental impact
of training and evaluation will be reported in the final manuscript.
